\documentclass[reqno, answers]{exam}

\usepackage{amssymb, amsmath, amsthm, stmaryrd}
\usepackage{fullpage, parskip}
\usepackage{graphicx}
\usepackage{enumitem}

% Improved spacing
\usepackage[top=1in, bottom=0.8in, left=1in, right=1in]{geometry}


\title{\normalsize{
  JHED: \underline{\hspace{4cm}} \hspace{1em} Name: \underline{\hspace{9cm}}}
   }
\author{}
\date{}

% Custom spacing for better page utilization
\renewcommand{\baselinestretch}{0.95}

\begin{document}
\maketitle
\vspace{-2cm} % Reduce space after title

\begin{questions}


  \question[10]
  We want to solve the following non-linear system of equations using Newton's method:
  \begin{align*}
    x + y^2 & = 1 \\
    x^2 - y & = 0
  \end{align*}
  The initial guess is $(x_0, y_0) = (1, 0)$. Find the first iteration of Newton's method, i.e. find $(x_1, y_1)$.

  \begin{solution}
    Let $\mathbf{F}(x, y) = \begin{bmatrix} x + y^2 - 1 \\ x^2 - y \end{bmatrix}$. The Jacobian matrix is given by
    \[J(x, y) = \begin{bmatrix} 1 & 2y \\ 2x & -1 \end{bmatrix}\]
    At the initial guess $(1, 0)$, we have
    \begin{align*}
      \mathbf{F}(1, 0) & = \begin{bmatrix} 0 \\ 1 \end{bmatrix}          \\
      J(1, 0)          & = \begin{bmatrix} 1 & 0 \\ 2 & -1 \end{bmatrix} \\
      \implies
      J(1, 0)^{-1}     & = \begin{bmatrix} 1 & 0 \\ 2 & -1 \end{bmatrix}
    \end{align*}

    Thus, the first iteration of Newton's method is given by
    \begin{align*}
      \begin{bmatrix} x_1 \\ y_1 \end{bmatrix} & = \begin{bmatrix} 1 \\ 0 \end{bmatrix} - J(1, 0)^{-1} \mathbf{F}(1, 0)                                                      \\
                                               & = \begin{bmatrix} 1 \\ 0 \end{bmatrix} - \begin{bmatrix} 1 & 0 \\ 2 & -1 \end{bmatrix} \begin{bmatrix} 0 \\ 1 \end{bmatrix} \\
                                               & = \begin{bmatrix} 1 \\ 0 \end{bmatrix} - \begin{bmatrix} 0 \\ -1 \end{bmatrix}                                              \\
                                               & = \begin{bmatrix} 1 \\ 1 \end{bmatrix}
    \end{align*}
  \end{solution}




  \newpage

  \question[5] We want to approximate $\pi$ by applying fixed-point iteration to $f(x)=\sin(x)$. Propose a function $g(x)$ for which the iteration converges to $\pi$ and is guaranteed to converge. Justify your choice.

  \begin{solution}
    Choose
    \[
      g(x)=x+c\sin(x),\qquad 0<c<2.
    \]
    Then $g(\pi)=\pi$, so $\pi$ is a fixed point. Also,
    \begin{align*}
      g'(x)   & = 1+c\cos(x), \\
      g'(\pi) & = 1-c.
    \end{align*}
    For local convergence to $\pi$, we need $|g'(\pi)|<1$, i.e.
    \[
      |1-c|<1 \iff 0<c<2.
    \]
    So any $c\in(0,2)$ works; for example, $c=1$ gives $g(x)=x+\sin(x)$.
  \end{solution}
  \vfill

  \question[5] Apply Newton's method near $x=0$ to $f(x)=\sin(x)$. Derive the error recurrence, i.e., express $\epsilon_{n+1}$ in terms of $\epsilon_n$. Your answer cannot be ``the error is always 0.'' You may use
  \[
    \sin(x)=\sum_{k=0}^{\infty}\frac{(-1)^k x^{2k+1}}{(2k+1)!},
    \qquad
    \tan(x)=x+\frac{x^3}{3}+\frac{2x^5}{15}+ \dots.
  \]

  \begin{solution}
    For Newton's method,
    \[
      x_{n+1}=x_n-\frac{\sin(x_n)}{\cos(x_n)}=x_n-\tan(x_n).
    \]
    Near $0$, use $\tan(x)=x+\frac{x^3}{3}+O(x^5)$:
    \[
      x_{n+1}=x_n-\left(x_n+\frac{x_n^3}{3}+O(x_n^5)\right)
      =-\frac{x_n^3}{3}+O(x_n^5).
    \]
    Since the root is $x^*=0$, the error is $\epsilon_n=x_n$, so
    \[
      \epsilon_{n+1}=-\frac{1}{3}\epsilon_n^3+O(\epsilon_n^5).
    \]
  \end{solution}
  \vfill

\end{questions}
\end{document}